\chapter{Implementation}

\section{Module Structure}

The processor is implemented as a modular SystemVerilog design with the following modules:

\begin{itemize}
    \item \texttt{single\_cycle.sv}: Top-level module integrating all components
    \item \texttt{register\_file.sv}: 32x32-bit register file
    \item \texttt{alu.sv}: Arithmetic Logic Unit
    \item \texttt{control\_unit.sv}: Instruction decoder and control signal generator
    \item \texttt{imem.sv}: Instruction memory
    \item \texttt{dmem.sv}: Data memory
\end{itemize}

\section{Register File Implementation}

The register file implements 32 registers with asynchronous read and synchronous write:

\begin{lstlisting}[language=Verilog, caption=Register File Implementation]
module register_file (
    input  logic        clk,
    input  logic        we,
    input  logic [4:0]  addr_rs1,
    input  logic [4:0]  addr_rs2,
    input  logic [4:0]  addr_rd,
    input  logic [31:0] wdata,
    output logic [31:0] rdata_rs1,
    output logic [31:0] rdata_rs2
);

    logic [31:0] registers [0:31];

    // Register x0 is always zero
    assign rdata_rs1 = (addr_rs1 == 5'b0) ? 32'b0 : registers[addr_rs1];
    assign rdata_rs2 = (addr_rs2 == 5'b0) ? 32'b0 : registers[addr_rs2];

    always_ff @(posedge clk) begin
        if (we && addr_rd != 5'b0) begin
            registers[addr_rd] <= wdata;
        end
    end

endmodule
\end{lstlisting}

\subsection{Key Features}

\begin{itemize}
    \item Register x0 hardwired to zero using combinational logic
    \item Writes to x0 are ignored (check \texttt{addr\_rd != 5'b0})
    \item Asynchronous read (data available immediately)
    \item Synchronous write (on clock edge)
\end{itemize}

\section{ALU Implementation}

The ALU supports 10 operations as specified in the requirements:

\begin{lstlisting}[language=Verilog, caption=ALU Operations]
always_comb begin
    case (alu_op)
        4'b0000: alu_out = op_a + op_b;           // ADD
        4'b0001: alu_out = op_a - op_b;           // SUB
        4'b0010: alu_out = op_a & op_b;           // AND
        4'b0011: alu_out = op_a | op_b;           // OR
        4'b0100: alu_out = op_a ^ op_b;           // XOR
        4'b0101: alu_out = op_a << op_b[4:0];     // SLL
        4'b0110: alu_out = op_a >> op_b[4:0];     // SRL
        4'b0111: alu_out = $signed(op_a) >>> op_b[4:0]; // SRA
        4'b1000: alu_out = ($signed(op_a) < $signed(op_b)) ? 32'b1 : 32'b0; // SLT
        4'b1001: alu_out = (op_a < op_b) ? 32'b1 : 32'b0; // SLTU
        default: alu_out = 32'b0;
    endcase
end
\end{lstlisting}

\subsection{Design Considerations}

\begin{itemize}
    \item Shift operations use lower 5 bits of op\_b for shift amount
    \item Arithmetic right shift (\texttt{SRA}) uses \texttt{\$signed()} for sign extension
    \item Comparison operations return 1 if condition true, 0 otherwise
    \item All operations are synthesizable
\end{itemize}

\section{Control Unit Implementation}

The control unit decodes instructions and generates control signals. The implementation uses a case statement on the opcode field:

\begin{lstlisting}[language=Verilog, caption=Control Unit Structure]
always_comb begin
    // Default values
    reg_write = 1'b0;
    mem_write = 1'b0;
    mem_read  = 1'b0;
    // ... other defaults ...
    
    case (opcode)
        7'b0110011: begin // R-type
            reg_write = 1'b1;
            // Set alu_op based on funct3 and funct7
        end
        7'b0010011: begin // I-type (arithmetic/logical)
            reg_write = 1'b1;
            alu_src_b = 2'b01; // Use immediate
            // Set alu_op based on funct3
        end
        7'b0000011: begin // Load instructions
            reg_write = 1'b1;
            mem_read  = 1'b1;
            mem_to_reg = 2'b01;
            // Set mem_size based on funct3
        end
        // ... other instruction types ...
    endcase
end
\end{lstlisting}

\subsection{Special Cases}

\begin{itemize}
    \item \textbf{LUI}: \texttt{alu\_src\_a = 2'b11} (immediate upper), \texttt{alu\_src\_b = 2'b11} (zero), \texttt{alu\_op = ADD}
    \item \textbf{AUIPC}: \texttt{alu\_src\_a = 2'b01} (PC), \texttt{alu\_src\_b = 2'b01} (immediate), \texttt{alu\_op = ADD}
    \item \textbf{JALR}: Similar to JAL but uses rs1 + immediate instead of PC + immediate
    \item \textbf{Shift instructions}: Use immediate[4:0] for shift amount (not full 32-bit immediate)
\end{itemize}

\section{Instruction Memory Implementation}

The instruction memory is implemented as a byte-addressable array:

\begin{lstlisting}[language=Verilog, caption=Instruction Memory]
module imem (
    input  logic [31:0]  addr,
    output logic [31:0]  rdata
);

    parameter MEM_SIZE = 16384;
    parameter ADDR_WIDTH = $clog2(MEM_SIZE);

    logic [7:0] mem [0:MEM_SIZE-1];
    logic [ADDR_WIDTH-1:0] byte_addr;

    initial begin
        $readmemh("../02_test/isa.mem", mem);
    end

    assign byte_addr = addr[ADDR_WIDTH-1:0];
    assign rdata = {mem[byte_addr+3], mem[byte_addr+2], 
                     mem[byte_addr+1], mem[byte_addr]};

endmodule
\end{lstlisting}

\subsection{Byte Ordering}

Instructions are stored in little-endian format (LSB at lower address), matching RISC-V specification.

\section{Data Memory Implementation}

The data memory supports byte, half-word, and word accesses with proper sign extension:

\begin{lstlisting}[language=Verilog, caption=Data Memory Write]
always_ff @(posedge clk) begin
    if (we && word_addr < MEM_SIZE) begin
        case (mem_size)
            3'b000: mem[word_addr] <= wdata[7:0];        // SB
            3'b001: begin                                // SH
                mem[word_addr]   <= wdata[7:0];
                mem[word_addr+1] <= wdata[15:8];
            end
            3'b010: begin                                // SW
                mem[word_addr]   <= wdata[7:0];
                mem[word_addr+1] <= wdata[15:8];
                mem[word_addr+2] <= wdata[23:16];
                mem[word_addr+3] <= wdata[31:24];
            end
        endcase
    end
end
\end{lstlisting}

\subsection{Memory Read with Sign Extension}

\begin{lstlisting}[language=Verilog, caption=Data Memory Read]
always_comb begin
    case (mem_size)
        3'b000: rdata = {{24{mem[word_addr][7]}}, mem[word_addr]}; // LB
        3'b001: rdata = {{16{mem[word_addr+1][7]}}, 
                         mem[word_addr+1], mem[word_addr]}; // LH
        3'b010: rdata = {mem[word_addr+3], mem[word_addr+2],
                         mem[word_addr+1], mem[word_addr]}; // LW
        3'b100: rdata = {24'b0, mem[word_addr]}; // LBU
        3'b101: rdata = {16'b0, mem[word_addr+1], mem[word_addr]}; // LHU
    endcase
end
\end{lstlisting}

\section{Memory-Mapped I/O Implementation}

The I/O system is integrated into the Load-Store Unit:

\subsection{I/O Write Logic}

\begin{lstlisting}[language=Verilog, caption=I/O Write Implementation]
always_ff @(posedge clk) begin
    if (~i_reset) begin
        o_io_ledr <= 32'b0;
        o_io_ledg <= 32'b0;
        o_io_lcd  <= 32'b0;
        // Reset HEX displays
    end else if (mem_write && dmem_addr >= 32'h1000_0000) begin
        if (dmem_addr >= 32'h1000_0000 && dmem_addr < 32'h1000_1000) begin
            o_io_ledr <= dmem_wdata;
        end else if (dmem_addr >= 32'h1000_1000 && dmem_addr < 32'h1000_2000) begin
            o_io_ledg <= dmem_wdata;
        end
        // ... LCD and HEX handling ...
    end
end
\end{lstlisting}

\subsection{I/O Read Logic}

\begin{lstlisting}[language=Verilog, caption=I/O Read Implementation]
always_comb begin
    if (dmem_addr >= 32'h1000_0000) begin
        if (dmem_addr >= 32'h1001_0000 && dmem_addr < 32'h1001_1000 && mem_read) begin
            io_data = i_io_sw;
        end else if (dmem_addr >= 32'h1000_4000 && dmem_addr < 32'h1000_5000 && mem_read) begin
            io_data = o_io_lcd;
        end else if (dmem_addr >= 32'h1000_2000 && dmem_addr < 32'h1000_3000 && mem_read) begin
            io_data = {1'b0, o_io_hex3, 1'b0, o_io_hex2, 
                       1'b0, o_io_hex1, 1'b0, o_io_hex0};
        end
        // ... other I/O regions ...
    end else begin
        io_data = dmem_rdata;
    end
end
\end{lstlisting}

\section{Top-Level Module Integration}

The \texttt{single\_cycle} module connects all components:

\begin{itemize}
    \item Instantiates all sub-modules (IMEM, register file, ALU, control unit, DMEM)
    \item Implements immediate generation logic
    \item Implements branch condition evaluation
    \item Implements PC update logic
    \item Manages memory-mapped I/O access
    \item Provides debug outputs (\texttt{o\_pc\_debug}, \texttt{o\_insn\_vld})
\end{itemize}

\subsection{Key Signal Assignments}

\begin{itemize}
    \item \texttt{o\_insn\_vld = i\_reset}: Instruction valid when not reset
    \item \texttt{o\_pc\_debug = pc}: Debug PC output
    \item PC updates on clock edge: \texttt{pc <= pc\_next}
    \item Register writes occur on clock edge when \texttt{reg\_write} is enabled
    \item Memory writes occur on clock edge when \texttt{mem\_write} is enabled
\end{itemize}

\section{Implementation Challenges and Solutions}

\subsection{Challenge 1: LUI Instruction}

\textbf{Problem}: LUI must load immediate to upper 20 bits, not add to any register value.

\textbf{Solution}: Use \texttt{alu\_src\_a = 2'b11} (immediate upper) and \texttt{alu\_src\_b = 2'b11} (zero), then ADD operation results in just the immediate value.

\subsection{Challenge 2: Memory-Mapped I/O}

\textbf{Problem}: Need to route memory accesses between DMEM and I/O regions.

\textbf{Solution}: Implement address decode logic in LSU to determine access type and route accordingly.

\subsection{Challenge 3: Sign Extension for Loads}

\textbf{Problem}: LB and LH must sign-extend, while LBU and LHU must zero-extend.

\textbf{Solution}: Use different \texttt{mem\_size} codes (3'b000 for LB, 3'b100 for LBU) and handle sign extension in read logic.

\subsection{Challenge 4: Misaligned Address Handling}

\textbf{Problem}: Addresses for word access must be word-aligned (LSB bits = 00).

\textbf{Solution}: Truncate address to word boundary using \texttt{word\_addr = \{addr[ADDR\_WIDTH-1:2], 2'b00\}}.

